% Flächeninhalt und Volumen im Raum · Zusammengesetzt und Verhältnisse · Prüfungs-Fokus Abitur GK – bank/flaecheninhalt-und-volumen-im-raum/e4.jsonl (zusammenbau v0.6, Prüfungs-Fokus)
\einheitenkopf[e4]{Einheit 4 · Einheit 4}
\zweigzeile{Abi GK ×5}
\setcounter{aufgabe}{0}

% Anlauf (ohne Prüfkennung)
\begin{aufgabe}{Zusammengesetzt und Verhältnisse – Anlauf}
\begin{teile}
\teil Ein Haus: Quader mit aufgesetztem Satteldach. Wie fasst du sein Volumen? \\ \kreuz{direkt} \\ \kreuz{als Summe} \\ \kreuz{als Differenz}
\teil Ein Quader hat die Grundfläche $PQRU$ mit $P(1 | 1 | 0)$, $Q(7 | 1 | 0)$, $R(7 | 7 | 0)$, $U(1 | 7 | 0)$ und die Höhe 4. Die Ebene durch $Q$, $U$ und $T(1 | 1 | 8)$ zerlegt den Quader. Berechne das Volumen des Teilkörpers, der $P$ enthält, und erläutere den Ansatz. V = \leerfeld VE
\teil Ein Würfel mit der Kantenlänge 8 hat eine Ecke im Ursprung, seine Kanten liegen auf den positiven Achsen. Die Ebene durch die Würfelkante auf der $x$-Achse und den Mittelpunkt $M(8 | 8 | 4)$ einer senkrechten Kante teilt den Würfel. Volumen des kleineren Teilkörpers? V = \leerfeld VE
\end{teile}
\end{aufgabe}

\begin{aufgabe}{Zusammengesetzt und Verhältnisse – Prüfungsaufgaben}
\begin{teile}
\teil Ein Quader hat die Grundfläche $PQRU$ mit $P(2 | 3 | 0)$, $Q(8 | 3 | 0)$, $R(8 | 9 | 0)$, $U(2 | 9 | 0)$ und die Höhe 5. Die Ebene durch $Q$, $U$ und $T(2 | 3 | 10)$ zerlegt den Quader. Berechne das Volumen des Teilkörpers, der $P$ enthält, und erläutere den Ansatz. \hfill (A18) \\ V = \leerfeld VE
\teil Ein Sockel hat die Form eines Stumpfs einer geraden quadratischen Pyramide: Grundfläche mit der Seite 6\,m, Deckfläche mit der Seite 3\,m, Höhe 6\,m. Sein Volumen wird mit dem Term $\tfrac{1}{3} \cdot 36 \cdot 12 - \tfrac{1}{3} \cdot 9 \cdot 6$ berechnet. Begründe, dass der Term das Volumen liefert. \hfill (A22)
\teil Die Punkte $(9 | 0 | 0)$, $(0 | 9 | 0)$ und $(0 | 0 | 9)$ liegen in der Ebene $x + y + z = 9$; die Punkte $K$, $L$, $N$ liegen je auf einer Koordinatenachse und in der Ebene $3x + 3y + 3z = 12$. Das Volumen des Körpers zwischen den beiden Ebenen im ersten Oktanten wird mit $\tfrac{1}{3} \cdot \tfrac{1}{2} \cdot 9^2 \cdot 9 - \tfrac{1}{3} \cdot \tfrac{1}{2} \cdot k^2 \cdot k$ berechnet. Erläutere den Ansatz und gib $k$ an. \hfill (A23)
\teil Ein Quader und eine Pyramide haben dieselbe Grundfläche, die Spitze der Pyramide liegt in der Deckfläche des Quaders. Um wie viel Prozent ist das Volumen des Quaders größer als das der Pyramide? Rechne ohne konkrete Volumina. \hfill (A21) \\ \leerfeld \%
\teil Das Dreieck mit den Ecken $(0 | 0)$, $(8 | 0)$ und $(8 | 6)$ rotiert um die $x$-Achse. Oberflächeninhalt des entstehenden Körpers? \hfill (A18) \\ O = \leerfeld FE
\teil Eine zylindrische Vorratsdose mit dem Durchmesser 20\,cm fasst 10 Liter. Passt sie in ein Regalfach mit 30\,cm lichter Höhe? \hfill (A22)
\end{teile}
\end{aufgabe}

\begin{aufgabe}{Zusammengesetzt und Verhältnisse – Prüfungsaufgaben – weiter}
\begin{teile}
\teil Für $-1 < k < 7$ hat die Pyramide $OP_kQ_kR$ die Ecken $P_k(7 - k | 0 | 0)$, $Q_k(0 | 1 + k | 0)$ und $R(0 | 0 | 3)$. Für welches $k$ ist ihr Volumen am größten? \hfill (A23) \\ k = \leerfeld
\teil Für $-4 < k < 8$ hat die Pyramide $OP_kQ_kR$ die Ecken $P_k(4 + k | 0 | 0)$, $Q_k(0 | 8 - k | 0)$ und $R(0 | 0 | 9)$. Für welches $k$ ist ihr Volumen am größten? \hfill (A23) \\ k = \leerfeld
\teil Ein Quader hat eine quadratische Grundfläche mit der Seite 6 in der $x$-$y$-Ebene und die Höhe 8. Auf der senkrechten Achse durch seine Mitte liegen $S_k$ in der Höhe $4 - k$ und $Q_k$ in der Höhe $4 + k$ ($0 \le k < 4$). Die Pyramide über der Grundfläche mit der Spitze $S_k$ und die Pyramide über der Deckfläche mit der Spitze $Q_k$ haben zusammen 25\,\% des Quadervolumens. Wert von $k$? \hfill (A26) \\ k = \leerfeld
\teil Ein Quader hat eine rechteckige Grundfläche mit den Seiten 4 und 5 in der $x$-$y$-Ebene und die Höhe 10. Auf der senkrechten Achse durch seine Mitte liegen $S_k$ in der Höhe $5 - k$ und $Q_k$ in der Höhe $5 + k$ ($0 \le k < 5$). Die Pyramide über der Grundfläche mit der Spitze $S_k$ und die Pyramide über der Deckfläche mit der Spitze $Q_k$ haben zusammen 20\,\% des Quadervolumens. Wert von $k$? \hfill (A26) \\ k = \leerfeld
\teil Aus einem Quader mit der Grundfläche 12 und der Höhe 6 wird die Pyramide über der Grundfläche mit der Spitze $S_t$ in der Höhe $6 - 2t$ herausgenommen ($0 \le t \le 3$). Die Abbildung zeigt zwei Graphen, einer gibt das Volumen des Restkörpers in Abhängigkeit von $t$ an. Welcher? Begründe. \hfill (A23) \\ \kreuz{Graph I} \\ \kreuz{Graph II}

\begin{ksys}[xmin=0,xmax=4,ymin=0,ymax=80,xstep=1,ystep=10,xlabel=t,ylabel=V] \funktionab{48+8*\x}{II}{0}{3} \funktionab{48+8/3*\x^2}{I}{0}{3} \end{ksys}
\teil Aus einem Quader mit der Grundfläche 12 und der Höhe 9 wird die Pyramide über der Grundfläche mit der Spitze $S_t$ in der Höhe $9 - 3t$ herausgenommen ($0 \le t \le 3$). Die Abbildung zeigt zwei Graphen, einer gibt das Volumen des Restkörpers in Abhängigkeit von $t$ an. Welcher? Begründe. \hfill (A23) \\ \kreuz{Graph I} \\ \kreuz{Graph II}

\begin{ksys}[xmin=0,xmax=4,ymin=0,ymax=120,xstep=1,ystep=20,xlabel=t,ylabel=V] \funktionab{72+12*\x}{I}{0}{3} \funktionab{72+4*\x^2}{II}{0}{3} \end{ksys}
\end{teile}
\end{aufgabe}

\medskip
%% Merkkasten Einheit 4: 6 Zeilen, Vorgabe höchstens fünf (3.1) – ungekürzt
\uebersichtskasten{Zerlegen und ergänzen: das gesuchte Volumen als Summe (Quader plus Dachprisma, Prisma plus Pyramide) oder als Differenz (große Pyramide minus überstehende Teile) fassen; vorgelegte Terme Faktor für Faktor den Teilkörpern zuordnen. \\ Verhältnisse ohne Zahlen: die Volumenformeln beider Körper mit denselben Größen aufstellen und dividieren – die Größen kürzen sich; Anteile am Quader über Grundflächenanteil und den Faktor ein Drittel begründen. \\ Ähnlichkeit: ein zur Grundfläche paralleler Schnitt trennt eine ähnliche Teilpyramide ab – jede Länge trägt den Streckfaktor, das Volumen seinen Kubus. \\ Teilpyramide mit dem Streckfaktor (4 − t)/4: Volumenhalbierung heißt ((4 − t)/4)³ = 1/2. \\ Parameter: das Volumen als Term im Parameter aufstellen – das Maximum liegt am Parabelscheitel (Mitte der Nullstellen), ein linearer Verlauf gehört zum Graphen als Gerade. \\ Rotation elementargeometrisch: das um eine Kathete gedrehte rechtwinklige Dreieck liefert einen Kegel (Radius, Höhe, Mantellinie über Pythagoras), die halbe Drehung den halben Kegel – Elementarformeln, kein Integral.}
